\documentclass[12pt,a4paper]{report}

%% ─── Packages ────────────────────────────────────────────────────────────────
\usepackage[utf8]{inputenc}
\usepackage[T1]{fontenc}
\usepackage[french,english]{babel}
\usepackage{geometry}
\usepackage{graphicx}
\usepackage{booktabs}
\usepackage{array}
\usepackage{enumitem}
\usepackage{hyperref}
\usepackage{fancyhdr}
\usepackage{titlesec}
\usepackage{setspace}
\usepackage{parskip}
\usepackage{xcolor}
\usepackage{microtype}

%% ─── Page geometry ───────────────────────────────────────────────────────────
\geometry{
  top=2.5cm, bottom=2.5cm,
  left=3cm,  right=2.5cm
}

%% ─── Headers & footers ───────────────────────────────────────────────────────
\pagestyle{fancy}
\fancyhf{}
\fancyhead[L]{\small EV Charge Manager}
\fancyhead[R]{\small \nouppercase{\leftmark}}
\fancyfoot[C]{\thepage}
\renewcommand{\headrulewidth}{0.4pt}

%% ─── Chapter title formatting ────────────────────────────────────────────────
\titleformat{\chapter}[block]
  {\normalfont\Large\bfseries}
  {CHAPTER \thechapter :}{0.8em}{}
\titlespacing*{\chapter}{0pt}{20pt}{15pt}

%% ─── Hyperref setup ──────────────────────────────────────────────────────────
\hypersetup{
  colorlinks = true,
  linkcolor  = blue!60!black,
  urlcolor   = blue!70!black,
  citecolor  = blue!60!black,
  pdftitle   = {EV Charge Manager -- MFE Fodhil \& Chibani},
  pdfauthor  = {Yacine FODHIL, Abderrahmane CHIBANI}
}

%% ─── Misc ────────────────────────────────────────────────────────────────────
\setlength{\parindent}{0pt}
\setlength{\parskip}{6pt}
\onehalfspacing

%%%============================================================================
\begin{document}
\selectlanguage{french}

%% ─── TITLE PAGE ──────────────────────────────────────────────────────────────
\begin{titlepage}
  \centering

  \textbf{Institut de Formation d'Assurances et de Gestion}\\[4pt]
  \textit{Établissement Privé de Formation Supérieure}\\[4pt]
  Agréé par le Ministère de l'Enseignement Supérieur et de la Recherche Scientifique

  \vspace{3cm}

  {\LARGE\bfseries MÉMOIRE DE FIN D'ÉTUDES}\\[8pt]
  {\large Pour l'obtention du diplôme de}\\[10pt]
  {\LARGE\bfseries LICENCE}\\[6pt]
  {\large Filière : Informatique}\\[4pt]
  {\large Spécialité : Métiers de l'Informatique}

  \vspace{2cm}

  {\large\bfseries THÈME}\\[12pt]
  \begin{center}
    \textbf{Smart Electric Vehicle Charging Station Management Platform:}\\[6pt]
    {\Large\textit{``EV Charge Manager''}}
  \end{center}

  \vspace{3cm}

  \begin{minipage}[t]{0.45\textwidth}
    \textbf{Réalisé Par :}\\[6pt]
    \begin{itemize}[leftmargin=1.2em, itemsep=2pt]
      \item Yacine FODHIL
      \item Abderrahmane CHIBANI
    \end{itemize}
  \end{minipage}
  \hfill
  \begin{minipage}[t]{0.5\textwidth}
    \textbf{Encadré Par :}\\[6pt]
    \begin{itemize}[leftmargin=1.2em, itemsep=2pt]
      \item Mme Nadia BENAHMED-EL ALIA
      \item M. Amar BALLA
    \end{itemize}
  \end{minipage}

  \vfill

  {\large Année Universitaire : 2025\,--\,2026}
\end{titlepage}

%% ─── TABLE OF CONTENTS ───────────────────────────────────────────────────────
\tableofcontents
\newpage

%%%============================================================================
%%  CHAPTER 1
%%%============================================================================
\selectlanguage{english}

\chapter{ISO 15118 Standard --- Definition \& Application in the EV Charging Platform Project}

\section{What is ISO 15118?}

ISO 15118 is an international standard that defines the communication interface between an
Electric Vehicle (EV) and Electric Vehicle Supply Equipment (EVSE)~--- commonly known as a
charging station. It is officially titled \textit{``Road Vehicles~--- Vehicle to Grid Communication
Interface''} and is developed by the International Organization for Standardization (ISO).

The standard goes far beyond simple plug-and-charge detection. It establishes a full
bidirectional digital communication protocol~--- referred to as Vehicle-to-Grid (V2G)
communication~--- enabling EVs and charging infrastructure to exchange rich data in real time.
This includes authentication, charge parameter negotiation, energy management, and support for
grid services.

ISO 15118 is structured across multiple parts (Parts 1--20), covering physical layer
requirements, network and transport layers, application layer messaging, and security. The
version most relevant to practical implementations is \textbf{ISO 15118-2}, which defines the
application-layer message set used during a charging session.

% ──────────────────────────────────────────────────────────────────────────────
\section{Why ISO 15118 Matters for Smart EV Charging}

Before ISO 15118, EV charging relied on the IEC 61851-1 standard, which uses simple low-level
control pilot signaling. This approach only communicates two pieces of information: whether a
vehicle is connected, and what maximum current the charger can supply. This severely limits
smart charging capabilities.

ISO 15118 addresses these limitations by enabling:

\begin{itemize}
  \item \textbf{Plug \& Charge (PnC):} The EV is automatically identified and authorized using
        digital certificates, eliminating the need for RFID cards or apps.

  \item \textbf{Bidirectional Communication:} Both the EV and EVSE can send and receive
        detailed parameters such as State of Charge (SoC), target current/voltage, and power
        limits.

  \item \textbf{Demand Response:} The charging station can dynamically adjust power delivery
        based on grid conditions or renewable energy availability.

  \item \textbf{Vehicle-to-Grid (V2G):} In supported configurations, energy stored in the EV
        battery can be fed back to the grid, turning EVs into distributed storage assets.

  \item \textbf{Security:} TLS-based encrypted communication and mutual certificate
        authentication protect the session end-to-end.
\end{itemize}

% ──────────────────────────────────────────────────────────────────────────────
\section{Application in Our EV Charging Platform}

Our Electric Vehicle Charging Platform uses ISO 15118 as its core communication backbone.
Rather than relying on simple current-signaling (IEC 61851-1), the platform implements full
V2G digital communication, enabling intelligent charge management and demand response
capabilities.

% ──────────────────────────────────────────────────────────────────────────────
\section{System Architecture}

The platform is built around two microcomputer modules~--- one representing the EVSE side
(\textbf{EVSEpi}) and one representing the EV side (\textbf{EVpi})~--- each running on a
Raspberry Pi 3 Model B. Communication between the SECC and EVCC is established over a Wi-Fi
wireless network, leveraging the open-source \textbf{RISE V2G} Java library (ISO 15118-2
implementation) as the communication engine.

Each module contains four sub-components:

\begin{itemize}
  \item \textbf{Communication Controller (SECC / EVCC):} Handles the ISO 15118 V2G message
        exchange. Modified to support dynamic parameter updates every 5 seconds, overcoming
        the static-parameter limitation of the base RISE V2G library.

  \item \textbf{Parameter File:} Acts as a shared interface between the communication
        controller and the management code. Stores all charging parameters (current, voltage,
        SoC, flags) in a format compatible with the RISE V2G library.

  \item \textbf{Management Code:} Java-based finite state machine that controls the charging
        process, updates parameters, and communicates with the HMI and IoT platform.

  \item \textbf{Human-Machine Interface (HMI):} Allows the EV user and Charge Point Operator
        (CPO) to configure parameters and monitor the charging session in real time.
\end{itemize}

% ──────────────────────────────────────────────────────────────────────────────
\section{DC Charging Mode and CC-CV Strategy}

The platform operates in DC power transfer mode, chosen because DC messaging carries a richer
parameter set than AC, allowing finer charge management. The EV module emulates a virtual
Li-ion battery using the industry-standard \textbf{Constant-Current / Constant-Voltage (CC-CV)}
charging strategy:

\begin{itemize}
  \item \textbf{CC Phase (SoC 0--79\%):} The EVCC requests maximum current while voltage rises
        linearly.

  \item \textbf{CV Phase (SoC 80--100\%):} Voltage is held at its maximum while current steps
        down in decreasing increments until the battery is full.
\end{itemize}

% ──────────────────────────────────────────────────────────────────────────────
\section{Dynamic Parameter Exchange}

A key enhancement made to the RISE V2G library is the ability to exchange dynamic parameters
during an active charging session. The system updates SoC, target current, target voltage, and
remaining charge time \textbf{every 5 seconds} via the \texttt{CurrentDemand} V2G message loop.
This enables the EVSE to respond in real time to changes in available power~--- critical for
renewable energy integration.

% ──────────────────────────────────────────────────────────────────────────────
\section{Demand Response with Photovoltaic Integration}

One of the primary use cases for ISO 15118 in our platform is demand response driven by
photovoltaic (PV) generation. When PV output fluctuates, the EVSEpi management code adjusts
the EVSE current and voltage limits accordingly. These updated limits are communicated to the EV
via the \texttt{CurrentDemandRes} message, causing the vehicle to adapt its charge rate in near
real time.

Testing with a high-resolution PV dataset confirmed that a 5-second communication interval
between SECC and EVCC produces a mismatch of only \textbf{0.38\%} between generation power and
charging power~--- an excellent balance between responsiveness and network efficiency
(0.2 packets/second).

% ──────────────────────────────────────────────────────────────────────────────
\section{IoT Integration and Monitoring}

Charging data is forwarded to an \textbf{Emoncms} IoT platform every minute by the EVSEpi
module. This allows the Charge Point Operator to monitor current, voltage, SoC, energy consumed
(kWh), and charging cost (EUR) over time via live dashboards~--- providing full operational
visibility of the ISO 15118 session data.

% ──────────────────────────────────────────────────────────────────────────────
\section{Security}

The platform uses certificate-based authentication as the payment/identity method, as defined
in the ISO 15118 Plug \& Charge (PnC) flow. Certificates were already implemented in the RISE
V2G library and provide greater security compared to External Identification Means (EIM),
making them the appropriate choice for a production-oriented charging platform.

% ──────────────────────────────────────────────────────────────────────────────
\section{Summary}

\begin{table}[h!]
  \centering
  \caption{ISO 15118 platform summary}
  \label{tab:iso15118-summary}
  \renewcommand{\arraystretch}{1.3}
  \begin{tabular}{>{\bfseries}p{4cm} p{9cm}}
    \toprule
    Aspect          & Detail \\
    \midrule
    Standard        & ISO 15118-2 (Application Layer --- DC Conductive Charging) \\
    Library         & RISE V2G (open-source Java implementation) \\
    Hardware        & Raspberry Pi 3 Model B (EVSEpi + EVpi) \\
    Communication   & Wi-Fi (replacing standard powerline; same ISO 15118 messages) \\
    Transfer Mode   & DC (richer parameter set for smart charging) \\
    Auth Method     & Certificate-based (Plug \& Charge) \\
    Update Rate     & 5-second charging loops (CC-CV dynamic updates) \\
    Demand Response & PV-integrated; 0.38\% power mismatch at 5\,s interval \\
    Monitoring      & Emoncms IoT platform (1-minute telemetry) \\
    \bottomrule
  \end{tabular}
\end{table}

% ──────────────────────────────────────────────────────────────────────────────
\section{References}

\begin{enumerate}[label={[\arabic*]}]
  \item Santos, J.B.; Francisco, A.M.B.; Cabrita, C.; Monteiro, J.; Pacheco, A.; Cardoso, P.J.S.
        Development and Implementation of a Smart Charging System for Electric Vehicles Based on
        the ISO 15118 Standard. \textit{Energies} 2024, 17, 3045.
        \url{https://doi.org/10.3390/en17123045}

  \item ISO 15118; Road Vehicles --- Vehicle to Grid Communication Interface (Parts 1--20).
        International Organization for Standardization (ISO): Geneva, Switzerland, 2022.

  \item SwitchEV. \textit{RISE-V2G Open-Source Library}.
        \url{https://github.com/SwitchEV/RISE-V2G}

  \item International Electrotechnical Commission (IEC). IEC 61851-1: Electric Vehicle
        Conductive Charging System --- Part 1: General Requirements. 2017.

  \item International Energy Agency. \textit{Global EV Outlook 2024}.
        \url{https://www.iea.org/reports/global-ev-outlook-2024}

  \item Mültin, M. ISO 15118 as the Enabler of Vehicle-to-Grid Application. In
        \textit{Proceedings of the 2018 International Conference of Electrical and Electronic
        Technologies for Automotive}, Milan, Italy, 9--11 July 2018.

  \item Emoncms IoT Platform. \url{https://emoncms.org/}
\end{enumerate}

%%%============================================================================
%%  CHAPTER 2
%%%============================================================================
\chapter{OCPP 1.6 vs.\ OCPP 2.0 --- A Comprehensive Comparison}

The Open Charge Point Protocol (OCPP) is the backbone of interoperability in the electric
vehicle (EV) charging industry, enabling seamless communication between charging stations and
central management systems (CSMS). While OCPP 1.6 laid the foundation for modern EV charging
networks, OCPP 2.0 represents a significant leap forward, introducing advanced features to
address the growing complexity of EV infrastructure.

In this chapter we explore the key differences between OCPP 1.6 and OCPP 2.0, highlighting
how the latter supports scalability, advanced energy management, and enhanced security.

% ──────────────────────────────────────────────────────────────────────────────
\section{Device Management and Structure}

\textbf{OCPP 1.6:}
\begin{itemize}
  \item OCPP 1.6 defines connectors as independent charging outlets, each treated as a
        standalone unit.
  \item There is no explicit concept of Electric Vehicle Supply Equipment (EVSE), limiting its
        ability to represent complex configurations like multi-connector stations.
\end{itemize}

\textbf{OCPP 2.0:}
\begin{itemize}
  \item OCPP 2.0 introduces a hierarchical device model, defining EVSE as a logical unit
        capable of managing one charging session at a time, even if multiple connectors are
        available.
  \item This formalized structure improves scalability and enables operators to manage advanced
        setups like multi-connector EVSEs and satellite systems efficiently.
\end{itemize}

% ──────────────────────────────────────────────────────────────────────────────
\section{Smart Charging Capabilities}

\textbf{OCPP 1.6:}
\begin{itemize}
  \item Smart charging is limited to predefined charging profiles:
        \begin{itemize}
          \item \textit{TxProfile}: Specific to a single charging session.
          \item \textit{TxDefaultProfile}: Applies to all transactions on a connector.
          \item \textit{ChargingStationMaxProfile}: Sets an upper limit on the charging
                station's total power consumption.
        \end{itemize}
  \item Charging profiles in OCPP 1.6 are static and cannot adapt dynamically to real-time
        conditions, making it less flexible for complex energy management scenarios.
\end{itemize}

\textbf{OCPP 2.0:}
\begin{itemize}
  \item Smart charging in OCPP 2.0 is more dynamic, allowing real-time updates to charging
        profiles based on factors such as grid conditions, energy prices, or vehicle schedules.
  \item \textit{Recurring Profiles}: Enables daily or weekly schedules for predictable
        charging patterns.
  \item Supports integration with ISO 15118 for Plug \& Charge and Vehicle-to-Grid (V2G)
        functionalities, providing bidirectional energy flow and automatic session initiation.
\end{itemize}

% ──────────────────────────────────────────────────────────────────────────────
\section{Message Handling and Communication}

\textbf{OCPP 1.6:}
\begin{itemize}
  \item Uses separate messages for related actions, such as \texttt{StartTransaction},
        \texttt{StopTransaction}, and \texttt{MeterValues}.
  \item Message structures are simpler but less efficient for large-scale operations.
\end{itemize}

\textbf{OCPP 2.0:}
\begin{itemize}
  \item Consolidates transaction-related messages into the unified \texttt{TransactionEvent}
        message, reducing redundancy and simplifying communication.
  \item Introduces new messages like \texttt{GetCompositeSchedule} for aggregating and
        visualising energy needs across a charging station.
  \item Supports WebSocket compression for more efficient communication, particularly useful
        for high-traffic networks.
\end{itemize}

% ── placeholder for use-case diagram ─────────────────────────────────────────
\begin{figure}[h!]
  \centering
  % \includegraphics[width=0.8\textwidth]{figures/usecase_auth_start.png}
  \framebox[0.8\textwidth]{\rule{0pt}{5cm}\quad [Figure: Use-case Diagram --- Auth using start button]\quad}
  \caption{Use-case Diagram: Authentication using the start button}
  \label{fig:usecase-auth}
\end{figure}

% ──────────────────────────────────────────────────────────────────────────────
\section{Enhanced Security}

\textbf{OCPP 1.6:}
\begin{itemize}
  \item Basic security measures include Transport Layer Security (TLS) for encrypted
        communication.
  \item Authentication is limited to username-password combinations, which can be vulnerable
        if not properly configured.
\end{itemize}

\textbf{OCPP 2.0:}
\begin{itemize}
  \item Implements robust security profiles:
        \begin{itemize}
          \item \textit{Security Profile 1}: Basic unencrypted communication (not recommended).
          \item \textit{Security Profile 2}: Encrypted communication using TLS.
          \item \textit{Security Profile 3}: End-to-end encryption and mutual authentication
                using X.509 certificates.
        \end{itemize}
  \item Introduces certificate management features for secure firmware updates and
        authentication, ensuring compliance with modern cybersecurity standards.
\end{itemize}

% ──────────────────────────────────────────────────────────────────────────────
\section{Firmware Management}

\textbf{OCPP 1.6:}
\begin{itemize}
  \item Supports remote firmware updates but lacks advanced monitoring and validation tools.
\end{itemize}

\textbf{OCPP 2.0:}
\begin{itemize}
  \item Enhances firmware management with features like secure delivery via HTTPS and
        signature verification to prevent unauthorized installations.
  \item Includes detailed status updates through \texttt{FirmwareStatusNotification},
        providing operators with real-time insights into the update process.
\end{itemize}

% ──────────────────────────────────────────────────────────────────────────────
\section{Backward Compatibility}

\textbf{OCPP 1.6:}
\begin{itemize}
  \item Compatible with earlier versions like OCPP 1.5, ensuring ease of transition for
        existing systems.
\end{itemize}

\textbf{OCPP 2.0:}
\begin{itemize}
  \item Not backward compatible with OCPP 1.6 due to significant structural and functional
        changes. Migration requires updating both CSMS and charger firmware.
\end{itemize}

% ──────────────────────────────────────────────────────────────────────────────
\section{Use Cases and Scalability}

\textbf{OCPP 1.6:}
\begin{itemize}
  \item Suitable for smaller, less complex charging networks with basic operational needs.
  \item Limited support for advanced use cases like V2G or complex energy optimization.
\end{itemize}

\textbf{OCPP 2.0:}
\begin{itemize}
  \item Designed for large-scale, modern EV networks with advanced requirements like fleet
        management, renewable energy integration, and demand response programs.
  \item Scalable for multi-connector EVSEs and capable of managing sophisticated energy use
        cases.
\end{itemize}

% ──────────────────────────────────────────────────────────────────────────────
\section{Technical Comparison}

\begin{table}[h!]
  \centering
  \caption{OCPP 1.6 vs.\ OCPP 2.0 --- Technical feature comparison}
  \label{tab:ocpp-comparison}
  \renewcommand{\arraystretch}{1.35}
  \begin{tabular}{>{\bfseries}p{5cm} p{4.5cm} p{4.5cm}}
    \toprule
    Feature                  & OCPP 1.6                  & OCPP 2.0 \\
    \midrule
    Release                  & 2015                      & 2020 \\
    Adoption                 & Very high                 & Still growing \\
    Smart charging           & Static profiles           & Dynamic real-time control \\
    Security                 & Basic TLS                 & Advanced certificate-based \\
    ISO 15118 integration    & Limited                   & Native support \\
    Device model             & Simple                    & Hierarchical \\
    Complexity               & Low                       & High \\
    \bottomrule
  \end{tabular}
\end{table}

Key takeaways from the comparison:
\begin{itemize}
  \item OCPP 2.0 introduces advanced smart charging, stronger security, and better device
        management, but requires new hardware and a major migration effort.
  \item OCPP 1.6 remains the dominant protocol in use worldwide.
  \item Most charging stations and backend systems support OCPP 1.6 by default.
  \item OCPP 2.0 is still in early adoption stages.
\end{itemize}

% ──────────────────────────────────────────────────────────────────────────────
\section{Why Using OCPP 1.6 in Algeria: Limited EV Infrastructure}

The choice of OCPP 1.6 over OCPP 2.0 in the Algerian context is justified by several
practical and technical considerations detailed below.

\begin{itemize}
  \item Algeria is still in the early EV deployment stage.
  \item No large smart-grid integration yet.
  \item No V2G infrastructure.
\end{itemize}

$\Rightarrow$ Therefore, the advanced features of OCPP 2.0 are not yet necessary.

\subsection{Hardware Availability}
\begin{itemize}
  \item Most low-cost chargers available in developing markets support only OCPP 1.6.
  \item OCPP 2.0 chargers are more expensive and rare.
\end{itemize}
$\Rightarrow$ OCPP 1.6 is economically suitable.

\subsection{Simplicity for Academic Prototype}
\begin{itemize}
  \item OCPP 1.6 is easier to implement.
  \item More documentation is available.
  \item Many open-source simulators use OCPP 1.6.
\end{itemize}

\subsection{Network Constraints (Important for Algeria)}
\begin{itemize}
  \item OCPP 1.6 works well with low bandwidth.
  \item OCPP 2.0 requires continuous advanced communication.
\end{itemize}
$\Rightarrow$ OCPP 1.6 is better for unstable internet environments.

\subsection{Industry Compatibility}

Most deployed chargers worldwide still run OCPP 1.6, ensuring interoperability between vendors.

% ──────────────────────────────────────────────────────────────────────────────
\section{Conclusion}

In this project, OCPP 1.6 was selected instead of OCPP 2.0 due to its widespread adoption,
lower implementation complexity, and better compatibility with existing charging infrastructure.
While OCPP 2.0 introduces advanced features such as enhanced security and dynamic smart
charging, these functionalities are not required in the current Algerian context where EV
deployment is still limited. Furthermore, OCPP 1.6 is supported by most available charging
stations and offers sufficient capabilities for remote monitoring, transaction management, and
basic smart charging. Therefore, OCPP 1.6 represents a practical and cost-effective choice for
the proposed system.

% ──────────────────────────────────────────────────────────────────────────────
\section{References}

\begin{enumerate}[label={[\arabic*]}]
  \item Luxman Energy. \textit{OCPP 1.6 vs OCPP 2.0: A Detailed Comparison for EV Chargers}.
        \url{https://www.luxmanenergy.com/ocpp-1-6-vs-ocpp-2-0-a-detailed-comparison-for-ev-chargers/}

  \item Ampcontrol. \textit{OCPP 1.6 vs OCPP 2.0: A Comprehensive Comparison}.
        \url{https://www.ampcontrol.io/post/ocpp-1-6-vs-ocpp-2-0-a-comprehensive-comparison}
\end{enumerate}

\end{document}
